\documentclass[aspectratio=169, 12pt]{beamer}

% --- Language and Fonts ---
\usepackage{fontspec}
\usepackage[english]{babel}
\usepackage{amsmath, amssymb, bm}
% tikz
\usepackage{tikz}
\usetikzlibrary{shapes, arrows.meta, positioning, shadows, calc}

% --- Bibliography (BibTeX + natbib) ---
\usepackage{natbib}
\bibliographystyle{chicago} 

% =========================================
% ===== keybox ====================
% =========================================
\usepackage[most]{tcolorbox}
\newtcolorbox{keybox}[1]{
  enhanced,
  colback=headbg!4,
  colframe=headbg,
  coltitle=headfg,
  colbacktitle=headbg,
  title=#1,
  fonttitle=\bfseries,
  boxrule=0.8pt,
  arc=1.5mm,
  left=2mm,
  right=2mm,
  top=1mm,
  bottom=1mm,
  before skip=0.2cm,
  after skip=0.2cm
}

\usepackage{booktabs}

% =========================================
% ===== Custom Footline (Page Number) =====
% =========================================


\setbeamercolor{footline}{fg=headbg}
\setbeamerfont{footline}{size=\scriptsize}

\setbeamertemplate{footline}{%
  \hfill%
  \usebeamercolor[fg]{footline}%
  \usebeamerfont{footline}%
  \insertframenumber\,/\,\inserttotalframenumber\kern1em\vskip3pt%
}

\setbeamertemplate{navigation symbols}{}

\usetheme{default}
\usepackage{xcolor}

\definecolor{headbg}{RGB}{30, 60, 120}
\definecolor{headfg}{RGB}{255, 255, 255}
\definecolor{accent}{RGB}{200, 50, 50}

\setbeamercolor{headline}{bg=headbg, fg=headfg}
\setbeamercolor{section in head/foot}{bg=headbg, fg=headfg}
\setbeamercolor{subsection in head/foot}{bg=headbg!80, fg=headfg}
\setbeamercolor{frametitle}{bg=headbg, fg=headfg}

% Step indicator command
\newcommand{\stepindicator}[1]{%
  \hfill {\scriptsize \textit{\textcolor{headbg}{Step #1 / 6}}}%
}
\DeclareMathOperator*{\argmin}{arg\,min}
\newcommand{\indep}{\perp\!\!\!\perp}

\usepackage{hyperref}


% --- Title ---
\title[]{A First Course in Causal Inference Ch.~11}
\author[Kodai Minesaki]{Kodai Minesaki}
\institute[Faculty of Engineering, The University of Tokyo]{Faculty of Engineering, The University of Tokyo}
\date{\today}

\begin{document}

% --- Title Slide ---
\begin{frame}
\titlepage
\end{frame}


\begin{frame}

  {
    \large
    Ch.~11 The Central Role of the Propensity Score in Observational Studies for Causal Effects
  }
  
\end{frame}

\begin{frame}{Ch.~11 The Central Role of the Propensity Score in Observational Studies for Causal Effects}
  \begin{itemize}
    \item \citet{rosenbaum1983central} proposed the key concept \textit{propensity score}.
    \item Under the i.i.d. sampling assumption, we have four random variables associated with each unit: $\{X,Z,Y(1),Y(0)\}$. 
    \item We can factorize the joint distribution as
    \[
    \Pr\{X,Z,Y(1),Y(0)\} = \Pr\{X\}\cdot\Pr\{Y(1),Y(0)\mid X\}\cdot\Pr\{Z\mid X,Y(1),Y(0)\}.
    \]
    \item We do not want to model the covariates becasue they are background information happening before the treatment and outcome. 
    \item If we want to go beyond the outcome model, then we must focus on the treatment assignment mechanism, which leads to the definition of the propensity score.
  \end{itemize}
\end{frame}

\begin{frame}
  \begin{center}
    \begin{keybox}{Definition~11.1 (propensity score)}
     Define 
     \[
     e(X,Y(1),Y(0))=\Pr\{Z=1\mid X,Y(1),Y(0)\}
     \] 
     as the propensity score. Under strong ignorability, we have 
     \[
     e(X,Y(1),Y(0))=\Pr\{Z=1\mid X,Y(1),Y(0)\} = \Pr(Z=1\mid X),
     \] 
     so the propensity score reduces to 
     \[
     e(X)=\Pr\{Z=1\mid X\},
     \]
     the conditional probability of receriving the treatment given the observed covariates.
  \end{keybox}
  \end{center}
\end{frame}

\begin{frame}{11.1 The propensity score as a dimension reduction tool}
  \begin{keybox}{Theorem~11.1}
    $Z\indep \{Y(1),Y(0)\}\mid X\Rightarrow Z\indep\{Y(1),Y(0)\mid e(X)\}.$
  \end{keybox}
  \begin{itemize}
    \item Thm.~11.1 implies that conditioning on the propensity score $e(X)$ removes all confunding induced by covariates $X$.
    \item The original covariates $X$ can be high-dimensional but $e(X)$ is one-dimensional scalar variable bounded between $0$ and $1$. 
    \item Therefore, the propensity score reduces the dimension of the original covariates but still maintains the ignorability: a \textit{dimensional reduction tool}.
  \end{itemize}
\end{frame}

\begin{frame}{11.1 The propensity score as a dimension reduction tool}
  \footnotesize
  \begin{proof}
    We need to show 
    \[
    \Pr\{Z=1\mid Y(1),Y(0),e(X)\} = \Pr\{Z=1\mid e(X)\}.
    \]
    The left-hand side of the equation equals:
    \begin{align*}
      \quad\Pr\{Z=1\mid Y(1),Y(0),e(X)\} 
      &= \mathbb{E}[Z\mid Y(1),Y(0),e(X)] \\
      &=\mathbb{E}[\mathbb{E}[Z\mid Y(1),Y(0),e(X),X]\mid Y(1),Y(0),e(X)]\quad (\because \text{LIE.})\\
      &= \mathbb{E}[\underbrace{\mathbb{E}[Z\mid X]}_{=e(X)}\mid Y(1),Y(0),e(X)]= e(X).
    \end{align*}
    The right-hand side of the equation equals:
    \begin{align*}
      \Pr\{Z=1\mid e(X)\}
      &= \mathbb{E}[Z\mid e(X)]
      = \mathbb{E}[\mathbb{E}[Z\mid e(X),X]\mid e(X)]\quad (\because \text{LIE.})= \mathbb{E}[\underbrace{\mathbb{E}[Z\mid X]}_{=e(X)}\mid e(X)]=e(X)
    \end{align*}
  \end{proof}
\end{frame}

\begin{frame}{11.1 The propensity score as a dimension reduction tool}
  \begin{itemize}
    \item Thm.~11.1 motivaes \textit{propensity score stratification}:
    \[
    Z\indep\{Y(1),Y(0)\}\mid e(X)=e_k\quad (k=1,\dots,K),
    \]
    where the propensity score only takes $K$ posible values $\{e_1,\dots,e_K\}.$
    \item In general, the propensity score is not known and is not discrete. 
    \item We often fit a statistical model for $\Pr\{Z=1\mid X\}$ to obtain $\hat{e}(X)$ and we can discretize it to approximate the simple case above.
    \item For example, we can discretize $\hat{e}(X)$ by its $K$ quantiles to obtain $\hat{e}'(X)$: $\hat{e}'(X)=e_k$, the $k/K$-th quantile of $\hat{e}(X)$, if $\hat{e}(X)$ is between the $(k-1)/K$-th and $k/K$-th quantitles of the $\hat{e}(X)$'s.
  \end{itemize}
\end{frame}

\begin{frame}{11.1 The propensity score as a dimension reduction tool}
  \begin{itemize}
    \item Its rigorous quantification is still missing in the literature.
    \item Q1: How do we choose $K$?
    \begin{enumerate}
      \item If $K$ is too small, ignorability does not hold conditional on $\hat{e}'(X)$ even approximately.
      \item If $K$ is too large, then we do not have enough units within each stratum of the estimated propensity score and many strata have only treated or control units.
    \end{enumerate}
    \item Q2: How do we compute standard errors of the estimators based on propensity score stratification?
    \begin{enumerate}
      \item Reporting standard errors based on the SRE ignores the uncertainty in the estimated propensity scores.
      \item Some researchers employ the bootstrap method but the theory for the bootstrap is still unclear due to the discreteness of this estimator.
    \end{enumerate}
  \end{itemize}
\end{frame}

\begin{frame}{11.2 Propensity score weighting}
  
\end{frame}

\begin{frame}{11.2 Propensity score weighting}
  
\end{frame}

\begin{frame}{11.2 Propensity score weighting}
  
\end{frame}

\begin{frame}{11.2 Propensity score weighting}
  
\end{frame}

\begin{frame}{11.2 Propensity score weighting}
  
\end{frame}

\begin{frame}{11.2 Propensity score weighting}
  
\end{frame}

\begin{frame}{11.2 Propensity score weighting}
  
\end{frame}


% =========================

% =========================
\section{}
\begin{frame}[allowframebreaks, noframenumbering]{References}
  \beamerdefaultoverlayspecification{}
  \small
  \bibliography{../references}
\end{frame}

\end{document}